\chapter{Аналитическая часть}

\section{Постановка задачи}
В лабораторной работе №2 необходимо провести сравнительный анализ рекурсивного и нерекурсивного алгоритма. По варианту необходимо найти максимальное число в последовательности чисел.

\section{Рекурсия}
\textbf{Рекурсия} — это такой способ организации вычислений, при котором процедура или функция в ходе выполнения обращается сама к себе. Другими словами, рекурсия является методом определения функций (процедур), при котором определяемая функция применена в теле своего же собственного определения [1].

\textbf{Рекурсивный алгоритм} — алгоритм, определяемый через себя [2].

\section{Механизм выполнения рекурсивных вызовов}
 Каждый вызов подпрограммы в первую очередь создаёт в стеке стековый кадр (stack frame), или фрейм — блок данных, размещаемый на вершине программного стека. Помимо адреса возврата, он содержит, в частности, параметры и локальные переменные метода или ссылки на них, а также другую низкоуровневую информацию. Адрес возврата указывает на машинную команду, которой метод должен передать управление по завершении своих действий. По завершении метода его запись активации удаляется из программного стека, а освобождаемая память может использоваться для вызовов других методов. Таким образом, в любой момент программный стек содержит информацию только об "активных"  подпрограммах, то есть вызванных, но ещё не завершённых (стековые кадры называют также активными кадрами, или записями активации) [3]. Глубина рекурсии определяется количеством вложенных вызовов и напрямую влияет на объем используемой памяти.
 
 \section*{Вывод}
В аналитической части рассмотрены основные понятия, связанные с рекурсией. Описана постановка задачи, проанализированы особенности рекурсивного подхода. Подчёркнута важность учёта глубины рекурсии при оценке трудоёмкости алгоритмов.
 
\clearpage
